\documentclass[UTF8,heading=true,12pt,font=songti]{article}
\usepackage{ctex}
\usepackage{pset}% 论文格式设置文件
\usepackage{hyperref} % 超链接默认方框

%封面信息
\serialnumber{20215034000} %学号
\college{数学与统计学院} %学院
\major{信息与计算科学} %专业
\grade{2021级} %年级
\author{NAME} %作者姓名
\titlefirst{标题} %题目第一行
\titlesecond{Subtitle} %题目第二行
\advisor{NAME} %指导老师
\advisors{Prof.} %职称
\submityear{25} %提交年份
\submitmonth{2} %提交月份
\submitdate{5}  %提交日
\zhtitle{中文标题}   %中文标题
\entitle{Title}	%英文标题

\begin{document}
	\maketitle% 论文封面
	%\includepdf[pages=-]{fig/titlepage}% word打印标题页导入
	\newpage
	\tableofcontents% 生成目录
	\setcounter{page}{0}% 此为第零页
	\thispagestyle{empty}% 页眉页脚为空
	\newpage
	\begin{center}
		\zihao{3}\bfseries{\the\zhtitle}%中文标题
	\end{center}
	
	\begin{table}[H]
		\vspace{-1.5em}
		\centering
		\begin{spacing}{1.5}
			\zihao{5}
			\begin{tabular}{rl} %设置表格字体位置(右左)
				学生姓名: {\the\author} & 学号: {\the\serialnumber}\\ 
				{\the\college} & {\the\major}专业\\
				指导教师: {\the\advisor} & 职称: {\the\advisors}\\	
			\end{tabular}
		\end{spacing}
		\vspace{-1.5em}
	\end{table}
	
	\addcontentsline{toc}{section}{摘\quad 要}\tolerance=500 % 将摘要放进目录
	{\heiti \zihao{-4}\bfseries{摘\quad 要}}:{\kaishu\zihao{-4}\ 摘要内容} % 缩进
	
	
	{\heiti \zihao{-4}\bfseries{关键词}}:{\kaishu\zihao{-4}\ 内容}% 缩进
	
	\begin{center}
		\zihao{3}{\the\entitle}%英文标题
	\end{center}
	\addcontentsline{toc}{section}{\bfseries Abstract}\tolerance=500 % 将Abstract放进目录
	
	{\zihao{-4}\bfseries{Abstract}}:\ \zihao{-4} Ab
	
	{\zihao{-4}\bfseries{Keywords}}:\ \zihao{-4} Ab
	
	\section{诸论}
	\subsection{研究背景}
	正文内容
	$$
	f(z_0)=\frac{1}{2\pi i}\oint_{\Gamma}\frac{f(z)}{z-z_0}\dd z.
	$$
	
	\section{预备知识}
	\subsection{常用记号}
	正文内容
	\begin{table}[H]
		\centering\small
		\renewcommand{\arraystretch}{1.5} %设置表格行距
		\caption{符号说明}
		\begin{tabular}{clc} %设置表格字体位置(中左中)
			\hline
			\textbf{符号} & \textbf{说明} & \textbf{单位} \\ \hline
			$\eta$ & 定日镜的光学效率 & 无量纲\\
			\hline
		\end{tabular}
	\end{table}  
	引入参考文献 \cite{bazlov2024momentsrepresentationnumbers}
	\begin{theorem}[thm]
		这是定理
	\end{theorem}
	\begin{lemma}
		hi
	\end{lemma}
	\begin{definition}
		hi
	\end{definition}
	\qu zm
	
	\pf 证明
	
	\subsubsection{小小一节}
	\begin{figure}[H]
		\centering
		\subfloat[示意图]{\includegraphics[width = 0.5\linewidth]{figure/logo-r.eps}}
		\hfill
		\subfloat[示意图]{\includegraphics[width = 0.5\linewidth]{figure/logo-r.eps}}
		\hfill
		\subfloat[示意图]{\includegraphics[width = 0.5\linewidth]{figure/logo-r.eps}}
		\hfill
		\subfloat[示意图]{\includegraphics[width = 0.5\linewidth]{figure/logo-r.eps}}
		\caption{示意图}
		\label{fig:label2}
	\end{figure}  
	
	
	\bibliography{ref/ref}
	\addcontentsline{toc}{section}{参考文献}\tolerance=500 % 将参考文献放进目录
	
\end{document}
    
